\begin{center}
\Large{\textbf{Abstract}}
\end{center}
\vspace{0.5cm}

The rapid evolution of cyber threats demands sophisticated, automated, and unified forensic analysis tools. This project presents \textbf{Vigil-Core}, an advanced forensics platform designed to streamline malware and network analysis for security analysts. Vigil-Core integrates multiple forensic capabilities into a single, cohesive full-stack application, replacing disjointed terminal-based tools with a dynamic, manager-ready web dashboard.

The system features a robust Python-based backend that handles heavy forensic workloads. It encompasses two core modules: \texttt{cyart\_malware} for static and dynamic malware analysis (including YARA rule matching and payload extraction), and \texttt{NetForensicX} for comprehensive network packet analysis. The frontend is built using React and Vite, offering a highly responsive, low-latency UI (averaging 4.3ms latency) that visualizes complex artifacts, such as entropy graphs and protocol distributions, using Recharts.

A key contribution to this project involved extensive Linux compatibility testing, end-to-end pipeline validation, and the architectural design of the UI. The platform was stress-tested with large real-world samples, including an 8GB memory dump, utilizing chunked processing to ensure stability. Furthermore, fuzzy hashing (pydeep/ssdeep) was integrated alongside YARA rules to detect modified malware variants accurately. 

Vigil-Core successfully bridges the gap between raw forensic data and actionable intelligence, providing an integrated environment that enhances the efficiency of Security Operations Centers (SOCs).
